\scanpage{108}
\mspage{54}

structures. On va donc \deleted{\unsure{fixer un|fixer une}}\inserted{pour} un \deleted{\unsure{instant}} \unsure{instant} oublier l'ens. $\Pi$ de structures 2-polyédrales sur $S$, et \deleted{\unsure{y}}\deleted{\unsure{retenir|retrouver}} partir d'une situation où on a \uline{une} telle structure

\[
\begin{aligned}
 &(S,F_0),\qquad &&S\text{ ensemble non vide}\\
 &&&F_0\subset\unsure{\mathscr P|\mathcal P|\mathfrak P}_3(S)\\
 &&&\text{faisant de }S\text{ un 2-polyèdre}\\
 &&&\text{triangulé}\qquad,
\end{aligned}\tag*{(20)}
\]

par la suite, revenant à $(S,F)$, on supposera que $F_0=F_{i_0}$, où $i_0\in\Pi$ est un élément particulier de $\Pi=\Pi_0(\Gamma)$. \deleted{Néanmoins,}\inserted{Les axiomes} que je vais envisager pour cette structure seront \deleted{\unsure{désig}}\inserted{distingués} par l'indice $0$.

\[
\begin{aligned}
 (\deleted{\unsure{S}}\mathrm T_0\,1)\quad&\text{Le 2-polyèdre }(S,F_0)\text{ est \uline{régulier},}\\
 &\text{i.e le groupe }\inserted{G_0}\text{ de ses automorphismes}\\
 &\text{est transitif (donc simplement}\\
 &\text{transitif sur l'ens des repères}\\
 &\operatorname{Rep}(P_0)
\end{aligned}\tag*{(21)}
\]

où on désigne par $P_0$ la structure polyédrale

\[
P_0=(S,A_0,F_0,R_{1,0},R_{2,0},R_{3\,0})\tag*{(22)}
\]

déduite de (20), et où on interprète $\operatorname{Rep}(P_0)$ comme

\[
\operatorname{Rep}(P_0)=\{(\unsure{s|a},t,u)\in S^3\mid\{s,t,u\}\in F_0\}.\tag*{(23)}
\]
% archive sheet number: 107
